\documentclass[11pt]{article}
\usepackage{icmlko}
\begin{document}

\title{두 해 만에 처음 그린 학습 곡선 --- 행을 두 배로 하면 수집 하나를 통째로 얻는 것과 같다}
\author{월드모델 실험실}
\date{노트 506 · 2026년 8월 3일}
\maketitle

\begin{abstract}
노트 505 가 모형 갈래를 닫았다(남은 여지 $+0.0009$). 그러면 남은 물음은
\textbf{무엇을 더 모으나}이고, 그 앞에 답해야 할 것이 하나 있다 ---
\textbf{같은 종류 행을 더 모으면 오르나.} 두 해 동안 학습 곡선을 한 번도
안 그렸다.

학습 행을 도메인 안에서 무작위로 $25 / 50 / 75 / 100\%$ 로 줄여 쟀다
(짝 씨앗 $6$ · 유보 고정 · 도메인 비율 유지).

\begin{center}
\begin{tabular}{lrrrl}
\hline
비율 & 학습 행 & 판 & 앞 칸 대비 & 부호 \\
\hline
$25\%$ & $4{,}645$ & $+0.4237$ & --- & --- \\
$50\%$ & $9{,}275$ & $+0.4598$ & $+0.0361$ & $6/6$ \\
$75\%$ & $13{,}904$ & $+0.4741$ & $+0.0143$ & $6/6$ \\
$100\%$ & $18{,}540$ & $\mathbf{+0.4843}$ & $\mathbf{+0.0102}$ & $\mathbf{6/6}$ \\
\hline
\end{tabular}
\end{center}

\textbf{마지막 칸이 $+0.0102$ 다.} 여덟 노트를 들인 모형 사이클(노트
$497 \sim 504$, $+0.0042$)의 \textbf{$2.4$ 배}이고 노트 505 가 잰 모형
남은 여지($+0.0009$)의 \textbf{$11$ 배}다.

뒤 세 점에 로그 직선을 맞추면 \textbf{두 배당 $+0.0245$} 다.

\begin{quote}
\textbf{학습 행을 두 배로 하면 $+0.025$ --- 검색 시계열 수집 전체
($0.0283$)와 맞먹는다.}
\end{quote}
\end{abstract}

\begin{figure}[t]
\centering
\includegraphics[width=\textwidth]{figs/f1.pdf}
\caption{\textbf{가} --- 곡선과 외삽. \textbf{나} --- 칸별 증가분.
\textbf{다} --- 손잡이 넷.}
\end{figure}

\section{손잡이 넷을 나란히}

\begin{center}
\begin{tabular}{lrl}
\hline
손잡이 & 판 크기 & 상태 \\
\hline
모형 남은 여지(노트 505) & $+0.0009$ & \textbf{닫힘} \\
여덟 노트 모형 사이클(노트 $497\!\sim\!504$) & $+0.0042$ & 다 씀 \\
\textbf{학습 행 두 배} & $\mathbf{+0.0245}$ & \textbf{열려 있음} \\
검색 시계열 수집 전체(노트 490) & $0.0283$ & 이미 있음(덮음 $41.7\%$) \\
\hline
\end{tabular}
\end{center}

\textbf{같은 종류 행을 두 배로 모으는 것과 새 수집을 통째로 하나 얻는 것이
같은 급이다.} 그리고 둘 다 모형 쪽에 남은 것의 서른 배 안팎이다.

\textbf{정직한 단서} --- 이 곡선은 \emph{줄여서} 그렸다. ``행이 적으면
이만큼''이지 ``행을 더하면 이만큼''이 아니다. 새로 모으는 행은 시기도
분포도 다를 수 있고(노트 493 의 덮음 체제 변화가 그런 일), 그러면 기울기가
작아진다. \textbf{$+0.0245$ 는 위 어림이다.}

다만 방향은 확실하다 --- \textbf{포화가 아니다.} 증가분이 $0.0361 \to
0.0143 \to 0.0102$ 로 줄지만 마지막 칸도 두 해 동안의 모형 작업을 다
합친 것보다 크다.

\section{이 실험실이 두 해 동안 무엇을 했나}

\begin{center}
\begin{tabular}{lr}
\hline
& 판 \\
\hline
채택된 모형 변경 다섯(계절 · 특기 · 깊이 · 애니 · 세계애니) & 합 $\approx +0.011$ \\
같은 기간 학습 행 $4{,}645 \to 18{,}540$ 이 준 것 & $\mathbf{+0.0606}$ \\
\hline
\end{tabular}
\end{center}

\textbf{자료가 모형보다 다섯 배 이상 일했다.} 두 해 동안 노트 $500$ 개를
모형에 썼고, 그 사이 조용히 늘어난 행이 그 다섯 배를 냈다.

\section*{남은 것}

\begin{itemize}
\item \textbf{어느 도메인의 행을 늘려야 하나} --- 이 곡선은 전 도메인을
같은 비율로 줄였다. 도메인별 곡선을 그리면 \emph{어디를 늘릴지}가 나온다.
\item \textbf{새 도메인 대 기존 도메인} --- 같은 행 수를 새 도메인으로
넣는 것과 기존 도메인에 넣는 것이 같은가. 노트 505 가 ``모으기는 작은
도메인을 끌어올린다''고 했으니 새 도메인이 더 클 수 있다.
\item \textbf{$+0.0245$ 를 위 어림에서 실측으로} --- 실제로 행을 더한
적이 없다. 다음 수집분이 들어오면 그 자리에서 확인한다.
\end{itemize}

\end{document}
